\documentclass[tikz,border=10pt]{standalone}
\usepackage{fontspec}
\usepackage{xeCJK}
\IfFontExistsTF{Noto Sans TC}{\setmainfont{Noto Sans TC}\setCJKmainfont{Noto Sans TC}}{\setmainfont{Microsoft JhengHei}\setCJKmainfont{Microsoft JhengHei}}
\xeCJKsetup{CJKglue={\hskip 0.08em plus 0.02em}, CJKecglue={\hskip 0.11em plus 0.02em}}
\usetikzlibrary{arrows.meta,positioning,calc}
\definecolor{paper}{HTML}{FFFDF8}
\definecolor{navy}{HTML}{111827}
\definecolor{muted}{HTML}{4B5563}
\definecolor{line}{HTML}{CBD5E1}
\definecolor{blue}{HTML}{1D4ED8}
\definecolor{bluebg}{HTML}{DBEAFE}
\definecolor{green}{HTML}{047857}
\definecolor{greenbg}{HTML}{D1FAE5}
\definecolor{gold}{HTML}{B45309}
\definecolor{goldbg}{HTML}{FEF3C7}
\definecolor{red}{HTML}{B91C1C}
\definecolor{redbg}{HTML}{FEE2E2}
\tikzset{
  title/.style={font=\fontsize{15}{20}\selectfont,text=navy,align=center},
  note/.style={font=\fontsize{9.5}{13}\selectfont,text=muted,align=center},
  stepbox/.style={rounded corners=12pt,line width=1pt,align=center,inner xsep=13pt,inner ysep=10pt,font=\fontsize{10.3}{13.8}\selectfont,text=navy}
}
\begin{document}
\begin{tikzpicture}[x=1cm,y=1cm,>=Latex,line cap=round,line join=round]
\path[fill=paper] (0,0) rectangle (16,9);

\node[title,text width=13cm] at (8,8.14) {agent 的價值，在一輪一輪修正};
\node[note,text width=11.6cm] at (8,7.52) {每一輪只要留下可查的證據，下一輪就不是重新猜。};

\node[stepbox,fill=bluebg,draw=blue,text width=2.25cm] (goal) at (2.4,5.2) {收到目標};
\node[stepbox,fill=greenbg,draw=green,text width=2.25cm] (act) at (5.95,5.2) {呼叫工具};
\node[stepbox,fill=goldbg,draw=gold,text width=2.25cm] (record) at (9.5,5.2) {留下紀錄};
\node[stepbox,fill=redbg,draw=red,text width=2.25cm] (check) at (13.05,5.2) {檢查結果};

\draw[->,draw=navy,line width=1.05pt] (goal)--(act);
\draw[->,draw=navy,line width=1.05pt] (act)--(record);
\draw[->,draw=navy,line width=1.05pt] (record)--(check);
\draw[->,rounded corners=14pt,draw=navy,line width=1.05pt] (check.south) -- (13.05,3.05) -- (2.4,3.05) -- (goal.south);

\node[stepbox,fill=white,draw=line,text width=5.4cm] at (8,2.0) {錯誤、差異、理由\\變成下一次修正的材料};
\node[note,text width=11.4cm] at (8,.78) {真正可靠的 agent，不是一次做對，而是錯了以後知道怎麼留下線索。};
\end{tikzpicture}
\end{document}
